% ======================================================================
% boards/rpi5.tex — Raspberry Pi 5 (BCM2712, Cortex-A76) のマルチコア化
% 実測: 2026-07-19, HTTP /bench (192.168.3.101)
% ======================================================================

Raspberry Pi 5（BCM2712、Cortex-A76 $\times$ 4、AArch64）へも同一のワーカープール型
SMP を移植した。ここでも\textbf{Pi 5 固有の差分と実測値}のみを述べる。

\subsection{対象と前提}

BCM2712 は 4 つの Cortex-A76 を持つ。A76 は DynamIQ Shared Unit（DSU）を備え、
D-cache ON 時のコア間コヒーレンシ機構が A53/A72 の \texttt{SMPEN} とは異なる（第
\ref{sec:dcache-a76} 節）。本節の測定は本レポート共通の前提どおり\textbf{D-cache
OFF}で行い、ロックフリー調停を RAM 直行で成立させている。A76 は 3 世代で\textbf{単一
コアの絶対速度が最も速い}（例：\bench{fill} 1\,MB の 1 コアが約 2{,}170\,\textmu s で、
A72 の約 7{,}000\,\textmu s の 3 倍以上速い）。

\subsection{コア起動：PSCI}
\label{sec:rpi5-bringup}

Pi 5 はコア起動に\textbf{PSCI}（\texttt{CPU\_ON} のハイパーコール）を用いる（Pi 3 の
armstub mailbox、Pi 4 の spin table に対応する差分）。コア識別は \texttt{MPIDR} の
Aff1 フィールドで行う。起床後の MMU 構成統一は他ボードと同一で、
\texttt{cores\_online = 4} を実機で確認した。

\subsection{実測 1：計算律速}

\begin{table}[htbp]
\centering
\caption{Pi 5：計算律速のコア数スイープ（\route{/bench}、\texttt{cores=1..4}、全て \texttt{agree=yes}）}
\label{tab:rpi5-compute}
\small
\begin{tabular}{@{}lccccl@{}}
\toprule
ベンチ & 1 & 2 & 3 & 4 & 備考 \\
\midrule
\bench{dining n=5}      & \spd{1.00} & \spd{1.99} & \spd{2.99} & \spd{3.99} & 完全に線形（理想） \\
\bench{primes n=300000} & \spd{1.00} & \spd{1.61} & \spd{2.31} & \spd{3.01} & 連続分割で概ね均衡 \\
\bench{nqueens n=12}    & \spd{1.00} & \spd{1.45} & \spd{1.47} & \spd{1.5}\,--\spd{1.5} & 静的分割では頭打ち \\
\bottomrule
\end{tabular}
\end{table}

\noindent
\bench{dining} は A53・A72 と同じく理想の \spd{3.99}。\bench{nqueens} は 4 コアで
\spd{1.5} 前後（\spd{1.41}--\spd{1.54}）と、\textbf{A72（\spd{1.85}）よりさらに低い}。
これは第 \ref{sec:cross-fair} 項の知見と整合する：A76 は単一コアが速いぶん、静的分割の
負荷不均衡とプール往復の相対比重が大きくなり、\bench{nqueens} の静的分割はより早く
頭打ちになる。改善には動的負荷分散（Level 3）が要る。

\subsection{実測 2：メモリ律速は最も平坦}

\begin{table}[htbp]
\centering
\caption{Pi 5：メモリ書込律速（\bench{fill}、8 パス、\texttt{cores=4}、5 回中の最良値）}
\label{tab:rpi5-fill}
\small
\begin{tabular}{@{}rrrl@{}}
\toprule
語数 & 1 コア \textmu s & 4 コア \textmu s & 速度向上 \\
\midrule
4{,}096  & 34    & 31    & \spd{1.09} \\
16{,}384 & 187   & 183   & \spd{1.02} \\
65{,}536 & 536   & 518   & \spd{1.09} \\
262{,}144 (=1\,MB) & 2{,}169 & 2{,}043 & \spd{1.06} \\
\bottomrule
\end{tabular}
\end{table}

\noindent
Pi 5 の \bench{fill} は全域で \spd{1.0}--\spd{1.1} と、3 世代で\textbf{最も平坦}で
ある。A76 は 1 コアあたりのメモリ帯域が相対的に高く、\textbf{1 コアで既に RAM 帯域を
ほぼ飽和}させるため、コアを足しても増えない。なお \bench{fill} の小サイズ（$\le$1{,}024
語、サブ\textmu s）は約 8 秒周期の自動アクタ GC に撹乱されて外れ値が出るため、5 回
測定の最良値で代表させた——単発測定のばらつき（第 \ref{sec:cross-fair} 項でも触れた）
の一例である。

\subsection{小括}

Pi 5（A76）でも本レポートの中心結論はそのまま成立する：\textbf{計算律速は理想的に
スケール（\spd{3.99}）・メモリ律速はスケールしない（\spd{1.06}）}。世代固有の差は、
起動機構（PSCI）、単一コア絶対速度が最速であること、\bench{fill} が最も平坦である
こと、そして D-cache ON 時のコヒーレンシ機構が DSU である点（第 \ref{sec:dcache-a76}
節）に限られる。「並列化すべき対象を分ける境界」は 3 世代を通じて同一であった。
